\section{Przygotowanie do gry}
\subsection{Ogólna zasada}
Każdy z graczy powinien zaopatrzyć się w ołówek oraz jedną kartę cech dla każdej postaci, którą gra.
Żetony powinny zostać starannie porozdzielane oraz posegregowane według rodzajów.
Zaleca się, aby wszystkie zapisy na kartach cech dokonywane były ołówkiem, tak, aby łatwo je było można zetrzeć, gdyż w trakcie gry cechy postaci ulegają ciągłym zmianom.

Przygotowanie do gry musi przebiegać w następujący sposób:

\subsection{Wybór bohaterów}
\begin{enumerate}
    \item Postacie są podzielone na dwie grupy: bohaterowie oraz uczniowie.
        Na żetonach reprezentujących bohaterów znajduje się imię, zaś żetony reprezentujące uczniów określają tylko ich rasę.
        Żetony reprezentujące bohaterów powinny zostać odwrócone zadrukowaną stroną ku dołowi oraz wymieszane tak, aby wybó była całkowicie przypadkowy.
    \item Jeżeli gra tylko jeden gracz, wybiera on trzech bohaterów (przypadkowo!) oraz trzech uczniów (dowolnie).
    \item Jeżeli w grze uczestniczy więcej niż jedna osoba, kolejność w jakiej dokonywany jest wybór bohaterów ustalana jest przy pomocy rzutu kostką.
        Wybór bohaterów dokonywany jest kolejno, zgodnie z ruchem wskazówek zegara, przy czym jako pierwszy dokonuje wyboru (,,ciągnie'') ten gracz, który wyrzucił największą liczbę oczek.
    \item Gracze ,,ciągną'' kolejno, aż do rozdzielenia między siebie trzech żetonów bohaterów i trzech żetonów uczniów.
        Liczba postaci (to znaczy bohaterów i uczniów) w grze może być mniejsza niz 6 (sugerowane minimum wynosi 4), jednak nie powinna przekraczać tej liczby.
\end{enumerate}

\subsection{Określenie cech postaci}
\begin{enumerate}
    \item Gracze odszukują cechy swoich bohaterów w tabeli \textit{Charakterystyka bohaterów} i zapisują uzyskane z tej tabeli informacje na karcie cech danego bohatera.
    \item Gracze wybierają rasę swoich uczniów.
        Uczeń może być krasnoludem, człowiekiem lub elfem.
    \item W zależności od rasy, na karcie cech ucznia zapisywane są następujące informacje:
        TABELA
\end{enumerate}
